% archon:covered-by Sawin_Terminology

% This chapter spans reference paper §3 and §4 (Sawin,
% sawin-perverse-sheaves-char-p.tex), whose verbatim titles organize it below.
% It proves the main theorem: the stalk dimension $\dim \cH^{-i}(K)_x$ of
% a perverse $\bF_\ell$-sheaf is bounded by the $i$th polar multiplicity of
% $\CCyc(K)$ at $x$ (Theorem~\ref{thm:sawin_first_massey}).  §3 establishes an
% alternate, manifestly well-defined form of the polar multiplicity; §4 carries
% out the pencil-of-conics induction.

\section{Equivalences between definitions of the polar multiplicity}
% Mirrors reference paper §3 "Equivalences between definitions of the polar
% multiplicity".

\begin{definition}[Local intersection number]
  \label{def:sawin_local_intersection}
  \lean{MR4228499.localIntersectionNumber}
  % SOURCE: Sawin, §3, Definition (local intersection number).
  % SOURCE QUOTE: We define their intersection number locally at $x$ \[ (C_1,C_2)_{Y, x} \] to be the sum of the degrees of the refined intersection $C_1 \cdot C_2$ on all connected components of $C_1 \cap C_2$ contained in $f^{-1}(x)$.
  \textit{Source: Sawin, §3, Definition (local intersection number).}
  For $f:Y\to X$, $x\in X$, and cycles $C_1,C_2$ on $Y$ of total dimension
  $\dim Y$ with $C_1\cap C_2\cap f^{-1}(x)$ proper (and each component of
  $C_1\cap C_2$ either proper inside $f^{-1}(x)$ or disjoint from it), the local
  intersection number $(C_1,C_2)_{Y,x}$ is the sum of the degrees of the refined
  intersection $C_1\cdot C_2$ (Fulton) over the components of $C_1\cap C_2$
  inside $f^{-1}(x)$.
\end{definition}

\begin{lemma}[Independence of the polar intersection from the choice of $(Y,V)$]
  \label{lem:sawin_intersection_comparison}
  \lean{MR4228499.intersection_comparison}
  \uses{def:sawin_projectivization, def:sawin_local_intersection}
  % SOURCE: Sawin, §3, Lemma intersection-comparison.
  % SOURCE QUOTE: For any $(Y,V)$ such that the strict transforms of $\mathbb P(C)$ and $\mathbb P(V)$ in the blowup of $\mathbb P(T^* X)$ at the fiber over $X$ do not intersect inside the exceptional divisor, the contribution of the fiber over $x$ to the intersection number $\mathbb P(C) \cap \mathbb P(V)$ depends only on $i$ and is independent of $Y,V$.
  \textit{Source: Sawin, §3, Lemma intersection-comparison.}
  Let $X$ be smooth, $C$ a conical cycle on $T^*X$ of dimension $\dim X$, $x\in
  X$, and $0\le i<\dim X$.  For $Y\subset X$ smooth of dimension $\dim X-i$
  through $x$ and $V$ a rank-$(i+1)$ sub-bundle of $T^*X$ on $Y$ whose strict
  transforms of $\PP(C),\PP(V)$ do not meet in the exceptional divisor of the
  blowup of $\PP(T^*X)$ at the fibre over $x$, the contribution of the fibre over
  $x$ to $\PP(C)\cap\PP(V)$ depends only on $i$.  A general pair $((TY)_x,V_x)$
  satisfies the hypothesis, so such $(Y,V)$ exist.
\end{lemma}
\begin{proof}
  \uses{def:sawin_projectivization}
  Work locally and deform $(Y,V)$ to $(Y',V')$ through a connected family of
  local complete intersections (convex combination of defining equations).  The
  emptiness of the strict-transform intersection in the exceptional divisor is
  open, so the whole family can be taken to satisfy it; the intersection locus is
  closed and disjoint from $x$, hence avoids a neighborhood of $x$, so the
  contribution at $x$ is constant in the family.  For genericity: $\PP(C)$ has
  dimension $\dim X-1$, the exceptional fibre $\PP((TX)_x)\times\PP((T^*X)_x)$ has
  dimension $2\dim X-2$, and the strict transform of $\PP(V)$ is
  $\PP((TY)_x)\times\PP(V_x)$ of dimension $\dim X-1$; for general subspaces the
  intersection has expected dimension $-1$, i.e. is empty.
\end{proof}

\begin{definition}[Polar multiplicity, intersection form]
  \label{def:sawin_polar_multiplicity_2}
  \lean{MR4228499.polarMultiplicity'}
  \uses{def:sawin_projectivization, def:sawin_local_intersection}
  % SOURCE: Sawin, §3, Definition polar-multiplicity-2.
  % SOURCE QUOTE: Define the $i$th polar multiplicity of $C$ at $x$ to be the intersection number \[ ( \mathbb P(C) , \mathbb P(V))_{ \mathbb P(T^* X), x} \] ... For $i=\dim X$, define the $i$th polar multiplicity ... to be the multiplicity of the zero section in $C$.
  \textit{Source: Sawin, §3, Definition polar-multiplicity-2.}
  For $0\le i<\dim X$, with $Y$ a sufficiently general smooth codimension-$i$
  subvariety through $x$ and $V$ a sufficiently general rank-$(i+1)$ sub-bundle of
  $T^*X$ over $Y$, the $i$th polar multiplicity of $C$ at $x$ is the local
  intersection number $(\PP(C),\PP(V))_{\PP(T^*X),x}$.  For $i=\dim X$ it is the
  multiplicity of the zero-section in $C$.  Well-definedness is
  Lemma~\ref{lem:sawin_intersection_comparison}.
\end{definition}

\begin{lemma}[Equivalence of the two polar multiplicities]
  \label{lem:sawin_multiplicity_polar}
  \lean{MR4228499.polarMultiplicity_eq}
  \uses{def:sawin_polar_multiplicity, def:sawin_polar_multiplicity_2, lem:sawin_intersection_comparison}
  % SOURCE: Sawin, §3, Lemma multiplicity-polar.
  % SOURCE QUOTE: Definitions polar-multiplicity and polar-multiplicity-2 are equivalent.
  \textit{Source: Sawin, §3, Lemma multiplicity-polar.}
  Definition~\ref{def:sawin_polar_multiplicity} and
  Definition~\ref{def:sawin_polar_multiplicity_2} agree.
\end{lemma}
\begin{proof}
  \uses{def:sawin_polar_multiplicity, def:sawin_polar_multiplicity_2}
  The multiplicity of a cycle at a point is its local intersection number with a
  general smooth subscheme through the point.  The projection formula gives
  $(\pi_*(\PP(C)\cap\PP(V)),Y)_{X,x} = (\PP(C),\PP(V)\cap\pi^*Y)_{\PP(T^*X),x}$,
  and $\PP(V)\cap\pi^*Y$ is the projectivized restriction $V'$ of $V$ to $Y$.  A
  dimension count in the exceptional divisor $\PP((TX)_x)\times\PP((T^*X)_x)$
  shows that for $V_x$ general and $Y$ general depending on $V$, the strict
  transforms of $\PP(C)$ and $\PP(V')$ do not meet, so
  Definition~\ref{def:sawin_polar_multiplicity_2} applies and the two numbers
  coincide.
\end{proof}

\section{A bound for Betti numbers}
% Mirrors reference paper §4 "A bound for Betti numbers".

\subsection{Transversality and nearby-cycle formulas}

\begin{lemma}[Fibre inclusion is properly transversal]
  \label{lem:sawin_transversality_check}
  \lean{MR4228499.transversality_check}
  \uses{def:sawin_C_transversal_map, def:sawin_properly_transversal}
  % SOURCE: Sawin, §4, Lemma transversality-check.
  % SOURCE QUOTE: If $f$ is $C$-transversal and the fibers of the composition $C \to X \to Y$ have dimension $m$, then $i$ is properly $C$-transversal.
  \textit{Source: Sawin, §4, Lemma transversality-check.}
  Let $f:X\to Y$ be smooth of relative dimension $m$ between smooth varieties,
  $C\subseteq T^*X$ closed conical with all components of dimension $\dim X$, and
  $i:f^{-1}(y)\hookrightarrow X$ the fibre inclusion.  If $f$ is $C$-transversal
  and the fibres of $C\to X\to Y$ have dimension $m$, then $i$ is properly
  $C$-transversal.
\end{lemma}
\begin{proof}
  \uses{def:sawin_C_transversal_map, def:sawin_properly_transversal}
  $C$-transversality of $f$ means $C$ meets the image of $df$ (the one-forms
  pulled back from $Y$, i.e. those in $\ker di$) only along the zero-section, so
  $i$ is $C$-transversal.  Moreover $i^*C = y\times_Y C$ is a fibre of $C\to Y$,
  of dimension $\dim X-\dim Y = \dim f^{-1}(y)$, giving proper transversality.
\end{proof}

\begin{lemma}[Nearby-cycle characteristic-cycle formula]
  \label{lem:sawin_nearby_formula}
  \lean{MR4228499.nearby_formula}
  \uses{lem:sawin_transversality_check, def:sawin_characteristic_cycle, def:sawin_shriek_pullback, def:sawin_singular_support, thm:saito_cc_nearby, thm:saito_index}
  % SOURCE: Sawin, §4, Lemma nearby-formula.
  % SOURCE QUOTE: Then \[i^! CC'(K) = CC ( R \Psi_{f} K) \] where we view the nearby cycles relative to $f$ as a complex of sheaves on $f^{-1}(Y)$.
  \textit{Source: Sawin, §4, Lemma nearby-formula.}
  Let $X$ be smooth, $Y$ a smooth curve, $K\in D^b_c(X,\bF_\ell)$, and let
  $\CCyc'(K)$, $\SSing'(K)$ denote $\CCyc(K)$, $\SSing(K)$ with the cotangent
  space at $x$ removed.  If $f:X\to Y$ is smooth projective, $\SSing'(K)$-transversal,
  and the fibres of $\SSing'(K)$ over $Y$ have dimension $\dim X-1$, then for
  $Z=f^{-1}(f(x))$ with inclusion $i$,
  \[ i^!\CCyc'(K) = \CCyc(\RPsi_f K). \]
\end{lemma}
\begin{proof}
  \uses{lem:sawin_transversality_check, def:sawin_characteristic_cycle}
  Away from $x$, $\SSing'(K)$-transversality makes $f$ locally acyclic, so
  $\RPhi_f K=0$ and $\RPsi_f K=i^*K$; by Saito Theorem 7.6 and
  Lemma~\ref{lem:sawin_transversality_check}, $\CCyc(\RPsi_f K)=i^!\CCyc'(K)$
  there.  Hence the difference $i^!\CCyc'(K)-\CCyc(\RPsi_f K)$ is a multiple of
  the cotangent space at $x$; pairing with the zero-section and using the index
  formula (Saito Theorem 7.13) $(\CCyc(\RPsi_f K),Z)=\chi(f^{-1}(\eta),K)$ on one
  side, and projection-formula manipulations $i_*(di)^*Z = N^*Z$ (Fulton
  8.1.1(c)) together with the constancy of $(\CCyc'(K),N^*f^{-1}(y'))$ in a smooth
  family on the other, shows the two pairings agree, forcing the multiple to be
  zero.
\end{proof}

\begin{lemma}[A general pencil of conics is transversal]
  \label{lem:sawin_genericity}
  \lean{MR4228499.genericity_lemma}
  \uses{def:sawin_properly_transversal, def:sawin_h_transversal, def:sawin_C_transversal_map}
  % SOURCE: Sawin, §4, Lemma genericity-lemma.
  % SOURCE QUOTE: Then $p$ is properly $C$-transversal, $q$ is $p^{\circ} C$-transversal in a neighborhood of the unique conic in the pencil containing $x$, and the fibers of $p^\circ C$ over $\mathbb P^1$ are $\dim X-1$-dimensional in a neighborhood of the conic containing $x$.
  \textit{Source: Sawin, §4, Lemma genericity-lemma.}
  Let $X\subseteq\PP^n$ be smooth, $C\subseteq T^*X$ closed conical of dimension
  $\dim X$, and $x\in X$ a point whose cotangent space is not contained in $C$.
  For a general pencil of conic sections $\overline X\subseteq X\times\PP^1$ with
  projections $p:\overline X\to X$, $q:\overline X\to\PP^1$: $p$ is properly
  $C$-transversal, $q$ is $p^\circ C$-transversal near the conic through $x$, and
  the fibres of $p^\circ C$ over $\PP^1$ are $(\dim X-1)$-dimensional near that
  conic.
\end{lemma}
\begin{proof}
  \uses{def:sawin_properly_transversal, def:sawin_h_transversal, def:sawin_C_transversal_map}
  The base locus $Y$ of a generic pencil is smooth of codimension $2$ (Bertini),
  and $p$ is the blow-up of $X$ along $Y$.  Counting conditions on the pencil:
  for $(y,\omega)\in C$, ``$y\in Y$'' is codimension $2$ and ``$\omega$ transverse
  to $Y$'' is codimension $\dim X-2$, while the space of such $(y,\omega)$ up to
  scaling is $(\dim X-1)$-dimensional, so the bad locus is codimension $1$ ---
  non-generic --- giving $C$-transversality; the dimension count of $p^*C$ gives
  proper transversality.  For $q$: $p^\circ C$-transversality is open, and one
  checks $dq(\omega_0)\notin p^\circ C$ pointwise on the conic through $x$ via
  analogous codimension counts, using that the cotangent space at $x$ is not in
  $C$.  Finally each component of $p^\circ C$ has dimension $\dim X$, so its
  fibres over $\PP^1$ have dimension $\dim X-1$ once the finitely many points
  forced into a single conic are excluded from the neighborhood.
\end{proof}

\begin{lemma}[Pencil-of-conics nearby/vanishing formula]
  \label{lem:sawin_final_nearby_formula}
  \lean{MR4228499.final_nearby_formula}
  \uses{lem:sawin_nearby_formula, lem:sawin_genericity, def:sawin_characteristic_cycle, def:sawin_shriek_pullback, thm:saito_cc_pullback, lem:gabber_illusie}
  % SOURCE: Sawin, §4, Lemma final-nearby-formula.
  % SOURCE QUOTE: (1) $CC ( R \Psi_{q} p^*K) = i^! p^! CC'(K)$ (2) $R\Phi_q p^* K$ is supported at $x$. (3) $-\operatorname{dimtot} \left( R\Phi_q ( p^* K) \right)_x$ is the multiplicity of the cotangent space at $x$ in $CC(K)$. (4) If $K$ is perverse, then $\left( R\Phi_q ( p^* K) \right)_x$ is supported in degree $-1$.
  \textit{Source: Sawin, §4, Lemma final-nearby-formula.}
  With $\overline X$, $p$, $q$ a general pencil of conics, $\overline x=p^{-1}(x)$,
  $y=q(\overline x)$, and $i:q^{-1}(y)\hookrightarrow\overline X$:
  (1) $\CCyc(\RPsi_q p^*K)=i^!p^!\CCyc'(K)$;
  (2) $\RPhi_q p^*K$ is supported at $x$;
  (3) $-\dimtot(\RPhi_q p^*K)_x$ equals the multiplicity of the cotangent space at
  $x$ in $\CCyc(K)$;
  (4) if $K$ is perverse, $(\RPhi_q p^*K)_x$ is supported in degree $-1$.
\end{lemma}
\begin{proof}
  \uses{lem:sawin_nearby_formula, lem:sawin_genericity, def:sawin_characteristic_cycle, def:sawin_dimtot}
  (1) By Lemma~\ref{lem:sawin_genericity}, $p$ is properly $\SSing(K)$-transversal
  and étale at $x$, so $\CCyc(p^*K)=p^!\CCyc(K)$ (Saito Theorem 6.6); apply
  Lemma~\ref{lem:sawin_nearby_formula} to $p^*K$ and $q$.
  (2) $q$ is $\SSing(p^*K)$-transversal near $x$ off $x$, so $p^*K$ is locally
  $q$-acyclic there and $\RPhi_q p^*K$ is supported at $x$.
  (3) $x$ is at most isolated characteristic, so
  $-\dimtot(\RPhi_q p^*K)_x = (\CCyc(p^*K),(dq)^*\omega)_x$; the only component of
  $\SSing(p^*K)$ meeting $(dq)^*\omega$ is $N^*x$, with intersection $1$, giving
  the multiplicity of $N^*x$ in $\CCyc(p^*K)=\CCyc(K)$.
  (4) $p^*K[-1]$ is perverse near $x$ on the generic fibre, so by Gabber's theorem
  (Illusie, Corollary 4.6) $\RPhi_q(p^*K)[-1]$ is perverse; being supported at a
  point it sits in degree $0$, so the unshifted complex is in degree $-1$.
\end{proof}

\begin{lemma}[Polar multiplicity drops under a generic conic section]
  \label{lem:sawin_multiplicity_hypersurface}
  \lean{MR4228499.multiplicity_hypersurface}
  \uses{lem:sawin_intersection_comparison, def:sawin_polar_multiplicity_2, def:sawin_shriek_pullback, def:sawin_polar_multiplicity}
  % SOURCE: Sawin, §4, Lemma multiplicity-hypersurface.
  % SOURCE QUOTE: Then for $i>0$, the $i$th polar multiplicity of $C$ at $x$ equals the $i-1$st polar multiplicity of $\inhyper{C}$ at $x$, and for $i=0$, the $i$th polar multiplicity of $C$ at $X$ equals the multiplicity of the cotangent space at $x$ in $C$.
  \textit{Source: Sawin, §4, Lemma multiplicity-hypersurface.}
  Let $X\subseteq\PP^n$ be smooth, $x\in X$, $C$ a conical cycle on $T^*X$, $C'$ its
  part away from the cotangent space at $x$, $\inhyper X$ a generic conic section
  through $x$ with inclusion $j$, and $\inhyper C=-j^!C'$.  Then for $i>0$ the
  $i$th polar multiplicity of $C$ at $x$ equals the $(i-1)$st polar multiplicity
  of $\inhyper C$ at $x$, and for $i=0$ it equals the multiplicity of the
  cotangent space at $x$ in $C$.
\end{lemma}
\begin{proof}
  \uses{lem:sawin_intersection_comparison, def:sawin_polar_multiplicity_2}
  Split into $i=0$, $0<i<\dim X$, $i=\dim X$.  For $i=0$ a general section of
  $\PP(T^*X)$ meets $\PP(C)$ only in the fibre over $x$.  For $0<i<\dim X$, a
  projective-geometry diagram (the blow-up of $\PP(T^*X\times_X\inhyper X)$ along
  the conormal section) and Fulton's projection formula identify
  $(\PP(\inhyper C),\PP(\inhyper V))_{\PP(T^*\inhyper X),x}$ with
  $(\PP(C'),\PP(V))_{\PP(T^*X)}=(\PP(C),\PP(V))_{\PP(T^*X)}$, while a codimension
  count on the exceptional divisor verifies the strict-transform condition of
  Lemma~\ref{lem:sawin_intersection_comparison}.  For $i=\dim X$, a cycle whose
  hypersurface restriction (resp. cotangent projection) is the zero-section was
  already the zero-section.
\end{proof}

\subsection{The main bound and its vanishing corollary}

\begin{theorem}[Characteristic-$p$ Massey bound]
  \label{thm:sawin_first_massey}
  \lean{MR4228499.first_Massey}
  \uses{lem:sawin_final_nearby_formula, lem:sawin_multiplicity_hypersurface, def:sawin_polar_multiplicity, def:sawin_characteristic_cycle}
  % SOURCE: Sawin, §4, Theorem first-Massey (= Theorem first-Massey-intro).
  % SOURCE QUOTE: Then $\dim_{\mathbb F_\ell} \mathcal H^{-i}(K)_x$ is at most $i$th polar multiplicity of $CC(K)$ at $x$.
  \textit{Source: Sawin, §4, Theorem first-Massey.}
  Let $X$ be a smooth variety over a perfect field $k$, $\ell$ invertible in $k$,
  and $K$ a perverse sheaf of $\bF_\ell$-modules on $X$.  Then for every $x\in X$,
  \[ \dim_{\bF_\ell}\cH^{-i}(K)_x \ \le\ \gamma^i_{\CCyc(K)}(x). \]
\end{theorem}
\begin{proof}
  \uses{lem:sawin_final_nearby_formula, lem:sawin_multiplicity_hypersurface}
  Étale-locally we may take $X\subseteq\PP^n$ smooth projective.  Induct on $i$.
  From the distinguished triangle $p^*K\to\RPsi_q p^*K\to\RPhi_q p^*K$ take stalk
  cohomology at $x$.  Since $\RPsi_q p^*K[-1]$ is perverse, $\cH^0(\RPsi_q p^*K)_x=0$,
  so for $i=0$, $\dim\cH^0(K)_x\le\dim(\RPhi_q p^*K)^{-1}_x\le-\dimtot(\RPhi_q
  p^*K)_x$, the multiplicity of the cotangent space at $x$ in $\CCyc(K)$, which by
  Lemma~\ref{lem:sawin_multiplicity_hypersurface} is $\gamma^0_{\CCyc(K)}(x)$.
  For $i>0$, Lemma~\ref{lem:sawin_final_nearby_formula}(4) makes the connecting map
  vanish, so $\dim\cH^{-i}(K)_x\le\dim\cH^{-i}(\RPsi_q p^*K)_x$; the induction
  hypothesis bounds this by the $(i-1)$st polar multiplicity of
  $\CCyc(\RPsi_q p^*K[-1])$, which by
  Lemma~\ref{lem:sawin_final_nearby_formula}(1) and
  Lemma~\ref{lem:sawin_multiplicity_hypersurface} equals $\gamma^i_{\CCyc(K)}(x)$.
\end{proof}

\begin{corollary}[Vanishing of stalk cohomology]
  \label{cor:sawin_vanishing}
  \lean{MR4228499.vanishing}
  \uses{thm:sawin_first_massey, def:sawin_polar_multiplicity, def:sawin_singular_support}
  % SOURCE: Sawin, §1/§4, Corollary vanishing-intro.
  % SOURCE QUOTE: Then $\mathcal H^{-i} (K)_x$ vanishes for \[ -i> \dim SS(K)_x - \dim X \] where $(SS(K))_x$ is the fiber of the singular support of $K$ over $x$.
  \textit{Source: Sawin, Corollary vanishing-intro.}
  With $X$, $k$, $\ell$, $K$ as in Theorem~\ref{thm:sawin_first_massey},
  $\cH^{-i}(K)_x = 0$ whenever $-i > \dim\SSing(K)_x - \dim X$, where
  $\SSing(K)_x$ is the fibre of the singular support over $x$.
\end{corollary}
\begin{proof}
  \uses{thm:sawin_first_massey}
  By Theorem~\ref{thm:sawin_first_massey} it suffices that the $i$th polar
  multiplicity vanishes for $i<\dim X-\dim\SSing(K)_x$.  For a rank-$(i+1)$ bundle
  $V$, $\PP(V)$ has dimension $i$ in the fibre over $0$ while $\PP(\CCyc(K))
  \subseteq\PP(\SSing(K))$ has codimension $\dim X-\dim\SSing(K)_x$ there, so a
  general $V$ misses it and the polar multiplicity is zero.
\end{proof}
